% Source: MQ15a_v1.html
\documentclass[11pt]{article}
\usepackage[margin=1in]{geometry}
\usepackage{amsmath}
\usepackage{amssymb}
\usepackage{siunitx}
\usepackage{enumitem}

\title{MQ15a: Waves}
\author{}
\date{}

\begin{document}
\maketitle

\section*{MQ15a: Waves}

\begin{enumerate}[leftmargin=*,label=\textbf{Question \arabic*.}]
\item \hfill \textit{1 pts}

Medical imaging systems often use ultrasound (frequencies above the range of human hearing). Calculate the frequency in MHz of an ultrasound wave traveling through human tissue at 1,550 m/s with a wavelength of 0.92 mm.



\emph{f} = \rule{2cm}{0.4pt} MHz

\item \hfill \textit{2 pts}

Suppose a wave has frequency 25.0 Hz and wavelength 3.30 m. Answer the following questions:

\textbf{(a) }The speed of the wave is [Select]
 m/s.

\textbf{(b)} The angular frequency is [Select]
 rad/s.

\textbf{(c)} The period is [Select]
 s.

\textbf{(d)} The wave number is [Select]
 rad/m.

\item \hfill \textit{1 pts}

A wave equation is given by $\alpha\frac{\partial^2y}{\partial t^2}-\beta\frac{\partial^2y}{\partial x^2}\:=\:0$ , where $\alpha=$ 10 s$^{2}$ and $\beta=$ 750.7 m$^{2}$. Calculate the wave speed.



\emph{v } = \rule{2cm}{0.4pt} m/s

\item \hfill \textit{3 pts}

Suppose a wave function is given by \emph{y}(\emph{t},\emph{x}) = 3 cos(2\emph{t }- 0.6\emph{x} + 1.2), where all values are consistent with the MKS system. Then ...

\textbf{(a)} the amplitude is [Select]
 m;

\textbf{(b)}the angular frequency is [Select]
 rad/s;

\textbf{(c)} the wave number is [Select]
 rad/m;

\textbf{(d) }the phase shift is [Select]
 rad;

\textbf{(e) }the frequency is [Select]
 Hz;

\textbf{(f)} the period is [Select]
 s;

\textbf{(g)} the wavelength is [Select]
 m;

\textbf{(h)} the propagation speed of the wave is [Select]
 m/s;

\textbf{(i)} the transverse \textbf{speed} of the wave at (0,0) is [Select]
 m/s.

\item \hfill \textit{1 pts}

A wire of mass 11 g and length 2.1 m under a tension of 454 N has an amplitude of 0.36 cm. If the frequency is 139 Hz, calculate the average power in watts of a wave propagating on the wire.



\emph{P}$_{av}$ = \rule{2cm}{0.4pt} W

\end{enumerate}

\end{document}
